% Options for packages loaded elsewhere
% Options for packages loaded elsewhere
\PassOptionsToPackage{unicode}{hyperref}
\PassOptionsToPackage{hyphens}{url}
\PassOptionsToPackage{dvipsnames,svgnames,x11names}{xcolor}
%
\documentclass[
  spanish,
  a4paper,
  oneside]{scrbook}
\usepackage{xcolor}
\usepackage[top=25mm,bottom=25mm,left=30mm,right=20mm]{geometry}
\usepackage{amsmath,amssymb}
\setcounter{secnumdepth}{5}
\usepackage{iftex}
\ifPDFTeX
  \usepackage[T1]{fontenc}
  \usepackage[utf8]{inputenc}
  \usepackage{textcomp} % provide euro and other symbols
\else % if luatex or xetex
  \usepackage{unicode-math} % this also loads fontspec
  \defaultfontfeatures{Scale=MatchLowercase}
  \defaultfontfeatures[\rmfamily]{Ligatures=TeX,Scale=1}
\fi
\usepackage{lmodern}
\ifPDFTeX\else
  % xetex/luatex font selection
\fi
% Use upquote if available, for straight quotes in verbatim environments
\IfFileExists{upquote.sty}{\usepackage{upquote}}{}
\IfFileExists{microtype.sty}{% use microtype if available
  \usepackage[]{microtype}
  \UseMicrotypeSet[protrusion]{basicmath} % disable protrusion for tt fonts
}{}
\makeatletter
\@ifundefined{KOMAClassName}{% if non-KOMA class
  \IfFileExists{parskip.sty}{%
    \usepackage{parskip}
  }{% else
    \setlength{\parindent}{0pt}
    \setlength{\parskip}{6pt plus 2pt minus 1pt}}
}{% if KOMA class
  \KOMAoptions{parskip=half}}
\makeatother
% Make \paragraph and \subparagraph free-standing
\makeatletter
\ifx\paragraph\undefined\else
  \let\oldparagraph\paragraph
  \renewcommand{\paragraph}{
    \@ifstar
      \xxxParagraphStar
      \xxxParagraphNoStar
  }
  \newcommand{\xxxParagraphStar}[1]{\oldparagraph*{#1}\mbox{}}
  \newcommand{\xxxParagraphNoStar}[1]{\oldparagraph{#1}\mbox{}}
\fi
\ifx\subparagraph\undefined\else
  \let\oldsubparagraph\subparagraph
  \renewcommand{\subparagraph}{
    \@ifstar
      \xxxSubParagraphStar
      \xxxSubParagraphNoStar
  }
  \newcommand{\xxxSubParagraphStar}[1]{\oldsubparagraph*{#1}\mbox{}}
  \newcommand{\xxxSubParagraphNoStar}[1]{\oldsubparagraph{#1}\mbox{}}
\fi
\makeatother


\usepackage{longtable,booktabs,array}
\newcounter{none} % for unnumbered tables
\usepackage{calc} % for calculating minipage widths
% Correct order of tables after \paragraph or \subparagraph
\usepackage{etoolbox}
\makeatletter
\patchcmd\longtable{\par}{\if@noskipsec\mbox{}\fi\par}{}{}
\makeatother
% Allow footnotes in longtable head/foot
\IfFileExists{footnotehyper.sty}{\usepackage{footnotehyper}}{\usepackage{footnote}}
\makesavenoteenv{longtable}
\usepackage{graphicx}
\makeatletter
\newsavebox\pandoc@box
\newcommand*\pandocbounded[1]{% scales image to fit in text height/width
  \sbox\pandoc@box{#1}%
  \Gscale@div\@tempa{\textheight}{\dimexpr\ht\pandoc@box+\dp\pandoc@box\relax}%
  \Gscale@div\@tempb{\linewidth}{\wd\pandoc@box}%
  \ifdim\@tempb\p@<\@tempa\p@\let\@tempa\@tempb\fi% select the smaller of both
  \ifdim\@tempa\p@<\p@\scalebox{\@tempa}{\usebox\pandoc@box}%
  \else\usebox{\pandoc@box}%
  \fi%
}
% Set default figure placement to htbp
\def\fps@figure{htbp}
\makeatother


% definitions for citeproc citations
\NewDocumentCommand\citeproctext{}{}
\NewDocumentCommand\citeproc{mm}{%
  \begingroup\def\citeproctext{#2}\cite{#1}\endgroup}
\makeatletter
 % allow citations to break across lines
 \let\@cite@ofmt\@firstofone
 % avoid brackets around text for \cite:
 \def\@biblabel#1{}
 \def\@cite#1#2{{#1\if@tempswa , #2\fi}}
\makeatother
\newlength{\cslhangindent}
\setlength{\cslhangindent}{1.5em}
\newlength{\csllabelwidth}
\setlength{\csllabelwidth}{3em}
\newenvironment{CSLReferences}[2] % #1 hanging-indent, #2 entry-spacing
 {\begin{list}{}{%
  \setlength{\itemindent}{0pt}
  \setlength{\leftmargin}{0pt}
  \setlength{\parsep}{0pt}
  % turn on hanging indent if param 1 is 1
  \ifodd #1
   \setlength{\leftmargin}{\cslhangindent}
   \setlength{\itemindent}{-1\cslhangindent}
  \fi
  % set entry spacing
  \setlength{\itemsep}{#2\baselineskip}}}
 {\end{list}}
\usepackage{calc}
\newcommand{\CSLBlock}[1]{\hfill\break\parbox[t]{\linewidth}{\strut\ignorespaces#1\strut}}
\newcommand{\CSLLeftMargin}[1]{\parbox[t]{\csllabelwidth}{\strut#1\strut}}
\newcommand{\CSLRightInline}[1]{\parbox[t]{\linewidth - \csllabelwidth}{\strut#1\strut}}
\newcommand{\CSLIndent}[1]{\hspace{\cslhangindent}#1}

\ifLuaTeX
\usepackage[bidi=basic,shorthands=off]{babel}
\else
\usepackage[bidi=default,shorthands=off]{babel}
\fi
\ifLuaTeX
  \usepackage{selnolig} % disable illegal ligatures
\fi


\setlength{\emergencystretch}{3em} % prevent overfull lines

\providecommand{\tightlist}{%
  \setlength{\itemsep}{0pt}\setlength{\parskip}{0pt}}



 


\AtBeginDocument{\hypersetup{linkcolor=black}}
% ---------- Paquetes tipograficos y microajustes ----------
\usepackage{microtype}        % mejor interletrado y justificado
\usepackage{csquotes}         % comillas tipograficas
\usepackage{iftex}
\ifPDFTeX\else
  % Seleccion de fuentes solo cuando XeTeX/LuaTeX esta disponible.
  \IfFontExistsTF{TeX Gyre Termes}{
    \setmainfont{TeX Gyre Termes}
  }{
    \IfFontExistsTF{Latin Modern Roman}{\setmainfont{Latin Modern Roman}}{}
  }
  \IfFontExistsTF{TeX Gyre Heros}{
    \setsansfont{TeX Gyre Heros}
  }{
    \IfFontExistsTF{Latin Modern Sans}{\setsansfont{Latin Modern Sans}}{}
  }
  \IfFontExistsTF{TeX Gyre Cursor}{
    \setmonofont{TeX Gyre Cursor}
  }{
    \IfFontExistsTF{Latin Modern Mono}{\setmonofont{Latin Modern Mono}}{}
  }
\fi
\usepackage{enumitem}         % listas compactas
\setlist{itemsep=.2em, topsep=.2em}

% ---------- KOMA-Script: estilo de titulos y espaciado ----------
\KOMAoptions{
  headings=big,
  parskip=half,
  fontsize=12pt,
  appendixprefix=true
}
\setlength{\parindent}{1.5em}
\setlength{\parskip}{0.6em}
\usepackage{indentfirst}
\usepackage{setspace}         % Control del interlineado del documento
\setstretch{1.5}

% ---------- Encabezados y pies (scrlayer-scrpage) ----------
\usepackage[automark,headsepline]{scrlayer-scrpage}
\clearpairofpagestyles
\automark[chapter]{chapter}
\ihead{\pagemark}
\ohead{\itshape\headmark}
\setheadsepline{0.4pt}
\renewcommand*{\chaptermarkformat}{}
\renewcommand*{\chapterpagestyle}{scrheadings}
\pagestyle{scrheadings}
\makeatletter
\let\ps@plain\ps@scrheadings
\makeatother
\cfoot{}

% ---------- Hipervinculos mas sobrios ----------
\usepackage{hyperref}
\hypersetup{
  colorlinks=true,
  linkcolor=black,
  citecolor=[rgb]{0.05,0.2,0.5},
  urlcolor=[rgb]{0.05,0.2,0.5},
  pdfauthor={\@author},
  pdftitle={\@title}
}

% ---------- Leyendas de figuras/tablas ----------
\usepackage[labelfont=bf,textfont=it]{caption}
\captionsetup{
  skip=8pt
}

% ---------- Tabla de contenidos ----------
\KOMAoptions{toc=graduated}
\RedeclareSectionCommand[tocnumwidth=3em]{chapter}
\RedeclareSectionCommand[tocindent=3.25em,tocnumwidth=2.8em]{section}
\RedeclareSectionCommand[tocindent=6.5em,tocnumwidth=2.5em]{subsection}
\makeatletter
\newcommand*{\tocdotfill}{\leavevmode\leaders\hbox to .6em{\hss.\hss}\hfill}
\makeatother
\RedeclareSectionCommand[toclinefill=\tocdotfill]{chapter}
\RedeclareSectionCommand[toclinefill=\tocdotfill]{section}
\RedeclareSectionCommand[toclinefill=\tocdotfill]{subsection}
\renewcommand*{\contentsname}{Tabla de contenido}
\setkomafont{chapterentry}{\normalfont}
\setkomafont{chapterentrypagenumber}{\normalfont}

% ---------- Viudas/Huerfanas y cortes de pagina ----------
\clubpenalty=10000
\widowpenalty=10000
\displaywidowpenalty=10000

% ---------- Entorno abstract para scrbook ----------
\providecommand{\abstractname}{Resumen}
\makeatletter
\@ifundefined{abstract}{
  \newenvironment{abstract}{
    \cleardoublepage
    \thispagestyle{plain}
    \null\vfill
    \begin{center}
      {\bfseries\Large \abstractname\par}
    \end{center}\vspace{1em}
    \begingroup
  }{
    \par\endgroup
    \vfill\null
    \cleardoublepage
  }
}{}
\makeatother

% ---------- Soporte de subtitulo desde YAML ----------
\makeatletter
\providecommand{\subtitle}[1]{\gdef\@subtitle{#1}}
\providecommand{\@subtitle}{}
\makeatother

\makeatletter
\newcommand{\facso@iflist}[4]{%
  \begingroup
    \def\addvspace##1{}%
    \def\numberline##1{##1}%
    \def\facso@target{#2}%
    \newcount\facso@entrycount
    \facso@entrycount=0\relax
    \long\def\contentsline##1##2##3{%
      \def\facso@this{##1}%
      \ifx\facso@this\facso@target
        \advance\facso@entrycount\@ne
      \fi
    }%
    \InputIfFileExists{\jobname.#1}{}{}%
  \endgroup
  \ifnum\facso@entrycount>0\relax
    #3%
  \else
    #4%
  \fi}
\renewcommand{\listoffigures}{%
  \facso@iflist{lof}{figure}{%
    \cleardoublepage
    \phantomsection
    \addchap*{Listado de Figuras}%
    \thispagestyle{scrheadings}%
    \markboth{Listado de Figuras}{Listado de Figuras}%
    \addcontentsline{toc}{chapter}{Listado de Figuras}%
    \@starttoc{lof}%
    \cleardoublepage
    \automark[chapter]{chapter}%
  }{}}
\renewcommand{\listoftables}{%
  \facso@iflist{lot}{table}{%
    \cleardoublepage
    \phantomsection
    \addchap*{Listado de Tablas}%
    \thispagestyle{scrheadings}%
    \markboth{Listado de Tablas}{Listado de Tablas}%
    \addcontentsline{toc}{chapter}{Listado de Tablas}%
    \@starttoc{lot}%
    \cleardoublepage
    \automark[chapter]{chapter}%
  }{}}
\makeatother

% --- Desactivar portada y abstract automaticos de Pandoc (PDF) ---
\AtBeginDocument{\let\maketitle\relax}
\renewenvironment{abstract}{}{}
\providecommand{\appendixname}{}
\providecommand{\appendixtocname}{}
\providecommand{\appendixpagename}{}
\renewcommand*{\appendixname}{Anexo}
\renewcommand*{\appendixtocname}{Anexos}
\renewcommand*{\appendixpagename}{Anexos}
\makeatletter
\@ifpackageloaded{bookmark}{}{\usepackage{bookmark}}
\makeatother
\makeatletter
\@ifpackageloaded{caption}{}{\usepackage{caption}}
\AtBeginDocument{%
\ifdefined\contentsname
  \renewcommand*\contentsname{Tabla de contenidos}
\else
  \newcommand\contentsname{Tabla de contenidos}
\fi
\ifdefined\listfigurename
  \renewcommand*\listfigurename{Lista de figuras}
\else
  \newcommand\listfigurename{Lista de figuras}
\fi
\ifdefined\listtablename
  \renewcommand*\listtablename{Lista de tablas}
\else
  \newcommand\listtablename{Lista de tablas}
\fi
\ifdefined\figurename
  \renewcommand*\figurename{Figura}
\else
  \newcommand\figurename{Figura}
\fi
\ifdefined\tablename
  \renewcommand*\tablename{Tabla}
\else
  \newcommand\tablename{Tabla}
\fi
}
\@ifpackageloaded{float}{}{\usepackage{float}}
\floatstyle{ruled}
\@ifundefined{c@chapter}{\newfloat{codelisting}{h}{lop}}{\newfloat{codelisting}{h}{lop}[chapter]}
\floatname{codelisting}{Listado}
\newcommand*\listoflistings{\listof{codelisting}{Listado de Listados}}
\makeatother
\makeatletter
\makeatother
\makeatletter
\@ifpackageloaded{caption}{}{\usepackage{caption}}
\@ifpackageloaded{subcaption}{}{\usepackage{subcaption}}
\makeatother
\usepackage{bookmark}
\IfFileExists{xurl.sty}{\usepackage{xurl}}{} % add URL line breaks if available
\urlstyle{same}
\hypersetup{
  pdftitle={¿Las mismas oportunidades para todos?},
  pdfauthor={Nicolás Outerbridge},
  pdflang={es},
  colorlinks=true,
  linkcolor={Maroon},
  filecolor={Maroon},
  citecolor={Blue},
  urlcolor={Blue},
  pdfcreator={LaTeX via pandoc}}


\title{¿Las mismas oportunidades para todos?}
\usepackage{etoolbox}
\makeatletter
\providecommand{\subtitle}[1]{% add subtitle to \maketitle
  \apptocmd{\@title}{\par {\large #1 \par}}{}{}
}
\makeatother
\subtitle{Creencias meritocráticas, experiencia escolar y percepción de
igualdad de oportunidades en estudiantes chilenos}
\author{Nicolás Outerbridge}
\date{22 de noviembre de 2026}
\begin{document}
\frontmatter
\maketitle

% --- title-pdf.tex personalizado para portada formal ---
\makeatletter
\providecommand{\subtitle}[1]{\gdef\@subtitle{#1}}
\providecommand{\@subtitle}{}
\providecommand{\frontmattercontext}{}
\providecommand{\advisorname}{}
\providecommand{\advisorlabel}{Profesor guia:}
\providecommand{\frontmatterlocation}{}
\InputIfFileExists{includes/cover-config.tex}{}{}

\newcommand{\PrintTitle}{%
  {\sffamily\bfseries\fontsize{24pt}{28pt}\selectfont \@title\par}%
}
\newcommand{\PrintSubtitle}{%
  \begingroup
  \edef\temp{\detokenize{\@subtitle}}%
  \ifx\temp\empty\relax
    % sin subtitulo
  \else
    {\sffamily\bfseries\large \@subtitle\par}%
  \fi
  \endgroup
}
\newcommand{\PrintAuthor}{%
  {\Large\bfseries \@author\par}%
}
\newcommand{\PrintDate}{%
  {\small \@date\par}%
}
\makeatother

\begin{titlepage}
\thispagestyle{empty}
\begin{center}
\vspace*{10mm}

% Logo o imagen institucional
\includegraphics[width=0.35\textwidth]{assets/cover.png}\par
\vspace{12mm}

% Titulo y subtitulo
\PrintTitle
\vspace{6mm}
\PrintSubtitle

\vspace{22mm}
\begingroup
\edef\temp{\detokenize{\frontmattercontext}}%
\ifx\temp\empty\relax
  % sin contexto adicional
\else
  {\normalsize \frontmattercontext\par}
\fi
\endgroup

\vspace{18mm}
\PrintAuthor

\vspace{16mm}
\begingroup
\edef\temp{\detokenize{\advisorname}}%
\ifx\temp\empty\relax
  % sin tutor
\else
  \rule{0.45\textwidth}{0.4pt}\par
  {\small \advisorlabel\ \advisorname\par}
\fi
\endgroup

\vfill
\begingroup
\edef\temp{\detokenize{\frontmatterlocation}}%
\ifx\temp\empty\relax
  % sin ubicacion
\else
  {\small \frontmatterlocation\par}
\fi
\endgroup
\PrintDate

\end{center}
\end{titlepage}

% Preliminares (numeros romanos) y estilo simple
\frontmatter
\pagestyle{scrheadings}

\setcounter{tocdepth}{2} % (o 1/3 segun prefieras)
\tableofcontents
\listoffigures
\listoftables


\mainmatter
\bookmarksetup{startatroot}

\chapter{title-block-style: none}\label{title-block-style-none}

p--- title: ``\,'' number-sections: false numbered: false format: html:
title-block-style: none toc: false pdf:

\small\textbf{Palabras clave:} igualdad de oportunidades, meritocracia,
privilegio, socialización escolar, educación superior, Chile \normalsize

\bookmarksetup{startatroot}

\chapter*{Resumen}\label{resumen}
\addcontentsline{toc}{chapter}{Resumen}

\markboth{Resumen}{Resumen}

La escuela promete que todos los estudiantes pueden llegar tan lejos
como su esfuerzo y su talento les permitan. Esa promesa descansa en un
supuesto anterior al mérito: que todos parten con las mismas
oportunidades. En Chile, sin embargo, el acceso a la educación superior
está fuertemente estratificado según el tipo de colegio de origen, y las
universidades más selectivas reclutan mayoritariamente a egresados de
colegios particulares pagados. Este estudio examina cómo perciben la
igualdad de oportunidades los estudiantes de básica y media, y qué
factores se asocian con esa percepción. Se analizan tres fuentes de
variación: la experiencia de justicia en la propia escuela; las
creencias meritocráticas de los estudiantes, distinguiendo percepciones
y preferencias sobre el mérito y el privilegio; y la posición en la
jerarquía escolar y el horizonte educativo. Para ello se utilizan datos
de las dos olas del panel Educación y Meritocracia (EDUMER), aplicado
entre 2023 y 2024 a estudiantes de básica y media de nueve escuelas
chilenas (N = 937 estudiantes y 1.576 observaciones).

\bookmarksetup{startatroot}

\chapter*{Prefacio}\label{prefacio}
\addcontentsline{toc}{chapter}{Prefacio}

\markboth{Prefacio}{Prefacio}

\mainmatter

\bookmarksetup{startatroot}

\chapter{Introducción}\label{introducciuxf3n}

Las instituciones educativas tienen la función explícita de transmitir
conocimientos y habilidades que preparan a los estudiantes para sus
futuras posiciones en la sociedad. Asimismo, transmiten normas, valores
y creencias que respaldan el logro educativo
(\citeproc{ref-mijs2016}{Mijs, 2016b}). Entre los principios centrales
destaca la promesa de que la escuela ofrece a todos la misma oportunidad
de demostrar su valía, y que las diferencias en los resultados reflejan
diferencias de esfuerzo y talento, no de origen
(\citeproc{ref-batruch_belief_2023}{Batruch et~al., 2023};
\citeproc{ref-bourdieu_reproduction_1990}{Bourdieu \& Passeron, 1990}).

Esta promesa reúne dos elementos distintos. El primero es la
\emph{meritocracia}, definida como la regla que distribuye recompensas
según esfuerzo y talento (\citeproc{ref-young1958}{Young, 1958}). El
segundo es la \emph{igualdad de oportunidades}, entendida como la
condición en la que todos compiten desde puntos de partida comparables,
de modo que el resultado dependa de las acciones individuales y no del
origen social. La igualdad de oportunidades es, así, una condición
previa de la meritocracia: solo cuando se cumple, recompensar el mérito
puede generar movilidad social
(\citeproc{ref-castillo_meritocracia_2019}{Castillo et~al., 2019}).
Ambos elementos, sin embargo, pueden separarse en la práctica, porque
una sociedad puede recompensar el esfuerzo y, al mismo tiempo,
distribuir desigualmente las oportunidades para demostrarlo. Mijs
(\citeproc{ref-mijs2016}{2016b}) lleva este argumento más lejos: como
las oportunidades para demostrar mérito dependen de factores externos a
la meritocracia, esta no solo es una promesa incumplida, sino también
incumplible.

El caso chileno ilustra esta tensión. El sistema escolar se caracteriza
por una alta estratificación y la noción de mérito ocupa un lugar
central en el discurso educativo
(\citeproc{ref-castillo_socialization_2024}{Castillo, Salgado, et~al.,
2024}). De manera paralela, la educación superior experimentó una rápida
masificación bajo un sistema jerarquizado, en el que pocas instituciones
forman a las élites y otras absorben a la mayoría
(\citeproc{ref-brunner_lucha_2012}{Brunner, 2012};
\citeproc{ref-brunner_dinamicas_2016}{Brunner \& Ganga, 2016}). Las
universidades más selectivas reclutan principalmente a egresados de
colegios particulares pagados, que representan una fracción menor de la
matrícula escolar total (\citeproc{ref-brunner_lucha_2012}{Brunner,
2012}; \citeproc{ref-munoz_taxonomia_2013}{Muñoz \& Blanco, 2013}). Para
los estudiantes de educación básica o media, la cuestión de si todos
tienen las mismas oportunidades no es abstracta, sino que se concreta en
el horizonte de acceso a distintas instituciones al finalizar la
escuela.

La investigación sobre creencias distributivas en la población escolar
ha examinado principalmente cómo los estudiantes evalúan la justicia y
el mérito en la propia escuela. La escuela es una ``esfera de justicia''
en la que se distribuyen constantemente recompensas, como las
calificaciones, cuya justicia es evaluada por los estudiantes
(\citeproc{ref-resh2010}{Resh, 2010}). Estas evaluaciones tienen
consecuencias que exceden el aula: quienes perciben a sus profesores
como justos se involucran más en actividades extracurriculares y de
voluntariado comunitario (\citeproc{ref-resh_sense_2017}{Resh \&
Sabbagh, 2017}). La creencia en la meritocracia escolar, es decir, que
en la escuela el éxito es resultado del mérito individual, está
extendida pese a la evidencia de que factores ajenos al mérito, como la
clase social, explican el éxito escolar, y opera como una herramienta de
justificación del sistema, especialmente entre los estudiantes de menor
estatus (\citeproc{ref-wiederkehr_belief_2015}{Wiederkehr et~al.,
2015}). Promover esta creencia puede incluso ampliar la brecha de
rendimiento según el origen social
(\citeproc{ref-darnon_where_2018}{Darnon et~al., 2018}), y sus efectos
se extienden a la evaluación de la sociedad, pues se asocia con percibir
como menos injusta la desigualdad social
(\citeproc{ref-batruch_belief_2023}{Batruch et~al., 2023}).

En Chile, esta línea de investigación ha examinado si los estudiantes
creen que el esfuerzo es recompensado y si consideran que debería serlo
(\citeproc{ref-castillo_socialization_2024}{Castillo, Salgado, et~al.,
2024}; \citeproc{ref-laffert_changes}{Laffert et~al., 2026}). Sus
resultados muestran que las creencias meritocráticas son
multidimensionales: los estudiantes distinguen entre lo que perciben y
lo que prefieren, y entre los principios del mérito y del privilegio, de
modo que pueden adherir al mérito como ideal y, al mismo tiempo,
reconocer el peso de la riqueza familiar y los contactos
(\citeproc{ref-laffert_changes}{Laffert et~al., 2026}). La percepción de
igualdad de oportunidades, en cambio, rara vez se ha estudiado como
objeto propio en población escolar, y con frecuencia se ha confundido
con la creencia meritocrática: parte de la literatura trata la
meritocracia como intercambiable con el apoyo genérico a la igualdad de
oportunidades (\citeproc{ref-laffert_changes}{Laffert et~al., 2026}).
Además, la evidencia sobre cómo los estudiantes perciben la desigualdad
en el acceso a la educación superior proviene mayoritariamente de
quienes ya han ingresado a ella, especialmente estudiantes de
universidades de élite (\citeproc{ref-allouch_ethique_2016}{Allouch
et~al., 2016}; \citeproc{ref-binder_symbolically_2019}{Binder \& Abel,
2019}; \citeproc{ref-villalobos_percepciones_2022}{Villalobos et~al.,
2022}). Existe escaso conocimiento sobre la formación de esta percepción
en etapas previas, cuando los estudiantes aún se encuentran en la
escuela y su trayectoria futura permanece abierta.

El presente estudio aborda esta brecha. Su pregunta principal es en qué
medida la percepción de igualdad de oportunidades de los estudiantes
escolares chilenos se asocia con su experiencia de meritocracia en la
escuela y con sus creencias sobre el mérito y el privilegio. De manera
secundaria, se examina si esa percepción se relaciona con la posición de
los estudiantes en la jerarquía escolar y con su horizonte educativo, y
cómo cambia a lo largo de la trayectoria escolar. Para responder a estas
preguntas se utilizan datos de las dos olas del panel Educación y
Meritocracia (EDUMER), que sigue a estudiantes de básica y media de
nueve escuelas chilenas.

\bookmarksetup{startatroot}

\chapter{Antecedentes conceptuales y
empíricos}\label{antecedentes-conceptuales-y-empuxedricos}

\section{Igualdad de oportunidades y
meritocracia}\label{igualdad-de-oportunidades-y-meritocracia}

La investigación sobre creencias meritocráticas ha distinguido, en
primer lugar, entre una dimensión normativa, la \emph{preferencia} por
que el esfuerzo y el talento sean la base de la distribución de
recompensas, y una dimensión descriptiva, la \emph{percepción} de que la
sociedad efectivamente los recompensa
(\citeproc{ref-castillo_meritocracia_2019}{Castillo et~al., 2019};
\citeproc{ref-duru-bellat_whos_2012}{Duru-Bellat \& Tenret, 2012}). La
percepción se refiere a lo que es observado (``lo que es'') y la
preferencia a lo que las personas quisieran que ocurriera (``lo que
debería ser'') (\citeproc{ref-castillo_meritocracia_2019}{Castillo
et~al., 2019}). Ambas dimensiones constituyen constructos diferentes: en
población adulta chilena, su correlación es baja y no significativa, y
la preferencia es más alta, más consensual y resistente, de modo que
puede mantenerse aunque se perciba que la meritocracia no opera
(\citeproc{ref-castillo_meritocracia_2019}{Castillo et~al., 2019}).

Castillo et al. (\citeproc{ref-castillo_multidimensional_2023}{2023})
agregan una segunda distinción, entre los elementos meritocráticos
(esfuerzo y talento) y los elementos basados en el privilegio (riqueza
familiar y contactos), que se tratan como componentes distintos y no
como polos opuestos de un mismo continuo
(\citeproc{ref-laffert_changes}{Laffert et~al., 2026}). De la
combinación de ambas distinciones resultan cuatro dimensiones: la
\emph{percepción de meritocracia} (si el esfuerzo y la habilidad son
recompensados), la \emph{percepción de privilegio} (si el éxito se
vincula a los contactos y la riqueza familiar), la \emph{preferencia
meritocrática} (si quienes se esfuerzan más o tienen más talento
deberían obtener mayores recompensas) y la \emph{aceptación del
privilegio} (si es aceptable que quienes tienen mejores contactos o
padres ricos obtengan mayor éxito)
(\citeproc{ref-laffert_changes}{Laffert et~al., 2026}). Esta estructura
de cuatro factores se ha validado en población escolar chilena con los
mismos datos que utiliza este estudio, y los resultados muestran que la
adhesión al mérito coexiste con el reconocimiento del privilegio, en
lugar de desplazarlo (\citeproc{ref-laffert_changes}{Laffert et~al.,
2026}).

Los niveles de cada dimensión ilustran esta coexistencia. En la primera
ola del panel EDUMER, la mayoría de los estudiantes percibe que en Chile
les va mejor a quienes tienen buenos contactos (77\%) o padres ricos
(71\%), y en menor medida que se recompensa el talento (74\%) o el
esfuerzo (67\%). Al mismo tiempo, el 87\% considera que el esfuerzo
debería ser la base de las recompensas, mientras que la aceptación de
que a quienes tienen padres ricos les vaya mejor se divide casi por
mitades y la de los buenos contactos es mayoritaria (61\%)
(\citeproc{ref-laffert_changes}{Laffert et~al., 2026}).

Las cuatro dimensiones de la creencia meritocrática no son
independientes entre sí, sino que se relacionan de manera diferenciada.
Entre los estudiantes, la percepción de meritocracia y la percepción de
privilegio presentan una correlación latente negativa (\emph{r} =
−0,21): quienes perciben que en la sociedad se recompensa el mérito
atribuyen un menor peso a los contactos y a la riqueza familiar. Al
mismo tiempo, la percepción de privilegio se asocia positivamente con la
preferencia meritocrática (\emph{r} = 0,48, la correlación más alta
entre las cuatro dimensiones), lo que se ha interpretado como una
respuesta normativo-reactiva ante la injusticia percibida: reconocer
ventajas injustas puede intensificar la demanda de que las recompensas
sigan al mérito (\citeproc{ref-castillo_multidimensional_2023}{Castillo
et~al., 2023}; \citeproc{ref-laffert_changes}{Laffert et~al., 2026}).

La percepción de igualdad de oportunidades es también una creencia
descriptiva, pero se refiere a otro momento del proceso distributivo.
Mientras la percepción de meritocracia evalúa la regla que conecta el
desempeño con la recompensa, la percepción de igualdad de oportunidades
evalúa las condiciones de partida: si el origen familiar, el colegio o
los recursos del hogar determinan quién puede competir y con qué
ventajas (\citeproc{ref-mijs2016}{Mijs, 2016b}). Por ello, ambas
creencias pueden disociarse. Un estudiante puede creer que en Chile el
esfuerzo es recompensado y, al mismo tiempo, reconocer que no todos
parten desde el mismo lugar.

Parte de la literatura ha operacionalizado la creencia meritocrática
mediante el apoyo genérico a la igualdad de oportunidades
(\citeproc{ref-batruch_belief_2023}{Batruch et~al., 2023};
\citeproc{ref-darnon_where_2018}{Darnon et~al., 2018}), lo que dificulta
distinguir ambas creencias (\citeproc{ref-laffert_changes}{Laffert
et~al., 2026}). Incluso las mediciones de la preferencia meritocrática
han incorporado la igualdad de oportunidades como condición, por
ejemplo, al preguntar si está bien que algunas personas tengan más
dinero ``pero sólo si hay igualdad de oportunidades''
(\citeproc{ref-castillo_meritocracia_2019}{Castillo et~al., 2019}).
Separar ambas creencias tiene consecuencias sustantivas. En la
investigación sobre percepciones de desigualdad, distintos indicadores
perceptuales se asocian débilmente y capturan aspectos diferentes del
fenómeno, con determinantes distintos
(\citeproc{ref-castillo_percepcion_2012}{Castillo et~al., 2012}). Este
estudio se centra en la percepción de igualdad de oportunidades y trata
las cuatro dimensiones de la creencia meritocrática como creencias
relacionadas pero distintas, que pueden contribuir a explicarla.

\section{Creencias meritocráticas y percepción de igualdad de
oportunidades}\label{creencias-meritocruxe1ticas-y-percepciuxf3n-de-igualdad-de-oportunidades}

Las cuatro dimensiones de la creencia meritocrática (percepción de
meritocracia, percepción de privilegio, preferencia meritocrática y
aceptación del privilegio) no tienen por qué relacionarse en la misma
dirección con la percepción de igualdad de oportunidades. La
meritocracia opera a la vez como un juicio cognitivo sobre cómo funciona
la desigualdad y como un lente moral para evaluar si los resultados
desiguales son merecidos (\citeproc{ref-laffert_changes}{Laffert et~al.,
2026}), de modo que cada dimensión puede relacionarse de forma distinta
con la evaluación de las oportunidades.

La percepción de meritocracia cumpliría una función de legitimación.
Quienes perciben que la meritocracia funciona adecuadamente perciben
menos desigualdad económica, lo que Castillo et al.
(\citeproc{ref-castillo_meritocracia_2019}{2019}) interpretan como una
legitimación de las diferencias de ingreso, que se deberían a razones
más bien individuales que sistémicas. Una percepción alta de
meritocracia es más compatible con la afirmación del orden existente,
mientras que una percepción baja puede reflejar insatisfacción y una
demanda de mayor justicia (\citeproc{ref-laffert_changes}{Laffert
et~al., 2026}). Si las recompensas se perciben como resultado del
esfuerzo y el talento, es esperable que también se perciba que las
condiciones para competir están abiertas a todos.

La percepción de privilegio apunta en la dirección contraria. Percibir
que a quienes tienen padres ricos o buenos contactos les va mejor
equivale a reconocer que el origen familiar y las redes condicionan los
resultados, es decir, que las oportunidades para demostrar mérito
dependen de factores externos a él (\citeproc{ref-mijs2016}{Mijs,
2016b}). La percepción de privilegio debería, por tanto, asociarse con
una menor percepción de igualdad de oportunidades.

La preferencia meritocrática no implica una evaluación positiva de la
realidad. En población adulta, una mayor preferencia por la meritocracia
se asocia con una mayor percepción de desigualdad, lo que sugiere que la
adhesión normativa al mérito establece un estándar más exigente frente a
la realidad y una demanda por reducir las diferencias que se basan en
criterios distintos al mérito
(\citeproc{ref-castillo_meritocracia_2019}{Castillo et~al., 2019}). En
población escolar, la preferencia meritocrática se asocia con una mayor
percepción de privilegio (\citeproc{ref-laffert_changes}{Laffert et~al.,
2026}). Quienes prefieren con más fuerza que el esfuerzo y el talento
determinen las recompensas podrían, así, evaluar más críticamente si
todos parten desde el mismo lugar. Sin embargo, la preferencia por la
meritocracia también ha sido descrita como un mito legitimador, que
provee justificaciones morales y racionales para la desigualdad social
(\citeproc{ref-castillo_meritocracia_2019}{Castillo et~al., 2019}), lo
que apunta en la dirección contraria. Con estos antecedentes no es
posible anticipar el sentido de su relación con la percepción de
igualdad de oportunidades.

La aceptación del privilegio es la dimensión menos estudiada, y los
antecedentes también admiten dos lecturas. Por una parte, quienes
consideran legítimo que el origen familiar y los contactos favorezcan a
algunos tendrían menos razones para interpretar esas ventajas como una
desigualdad de oportunidades. Por otra, en población escolar la
aceptación del privilegio se relaciona positivamente con la percepción
de privilegio, lo que indica que reconocer mecanismos no meritocráticos
se vincula con aceptar en mayor medida su legitimidad
(\citeproc{ref-laffert_changes}{Laffert et~al., 2026}); quienes aceptan
el privilegio también reconocerían que no todos parten desde el mismo
lugar. Como esta aceptación puede operar como un juicio normativo
separado de la evaluación del mérito
(\citeproc{ref-laffert_changes}{Laffert et~al., 2026}), su relación con
la percepción de igualdad de oportunidades se examina de manera
exploratoria.

\section{La escuela como espacio de socialización de las
oportunidades}\label{la-escuela-como-espacio-de-socializaciuxf3n-de-las-oportunidades}

La escuela es la institución que de manera más explícita promete
igualdad de condiciones: todos reciben el mismo currículo, son evaluados
con las mismas pruebas y compiten por las mismas notas. Su estructura,
su discurso y sus prácticas de evaluación contribuyen a presentar el
éxito académico como resultado del talento y el esfuerzo individual, y a
transformar diferencias de origen en diferencias de desempeño que
aparecen como individuales
(\citeproc{ref-bourdieu_reproduction_1990}{Bourdieu \& Passeron, 1990};
\citeproc{ref-goudeau_hidden_2017}{Goudeau \& Croizet, 2017}). La
meritocracia escolar suele presentarse como una forma de combatir la
desigualdad, pero la escuela lleva a los estudiantes a atribuir su
desempeño a factores internos, como el esfuerzo y la habilidad
(\citeproc{ref-darnon_where_2018}{Darnon et~al., 2018}).

El ideal de la escuela como gran igualadora no se realiza en la práctica
(\citeproc{ref-resh2010}{Resh, 2010}). Las oportunidades para demostrar
mérito dependen del origen social (\citeproc{ref-mijs2016}{Mijs,
2016b}), y las prácticas de aula convierten ventajas de clase en
diferencias de rendimiento (\citeproc{ref-goudeau_hidden_2017}{Goudeau
\& Croizet, 2017}). La experiencia de justicia escolar tampoco es
homogénea: en Israel, cerca de la mitad de los estudiantes de secundaria
se considera perjudicada en sus calificaciones
(\citeproc{ref-resh2010}{Resh, 2010}). La forma en que la escuela
organiza a los estudiantes también moldea sus explicaciones: en sistemas
que separan a los estudiantes según rendimiento, estos tienden a
atribuir su fracaso a sí mismos
(\citeproc{ref-mijs_stratified_2016}{Mijs, 2016a}).

Lo que ocurre en la escuela parece extenderse más allá de ella. Entre
estudiantes israelíes de octavo y noveno grado, quienes perciben a sus
profesores como justos se involucran más en actividades
extracurriculares y de voluntariado comunitario
(\citeproc{ref-resh_sense_2017}{Resh \& Sabbagh, 2017}). Creer que la
escuela es meritocrática se asocia con percibir como menos injusta la
desigualdad social y con un menor apoyo a la redistribución, y esta
relación se mantiene al controlar por la creencia en la meritocracia
general (\citeproc{ref-batruch_belief_2023}{Batruch et~al., 2023}).
Batruch et al.~interpretan que creer que la escuela es justa implica
creer que también lo son las diferencias que produce la jerarquía
educativa. Si la escuela es la primera institución en que los
estudiantes experimentan una distribución sistemática de recompensas, la
evaluación de que esa distribución es justa podría extenderse a la
percepción de que la sociedad ofrece a todos las mismas oportunidades.
Este argumento no pone en cuestión que la escuela deba ser justa, sino
que llama la atención sobre una implicancia de esa justicia: si la
experiencia de una distribución justa en la escuela se generaliza a la
sociedad, la propia equidad escolar podría llevar a los estudiantes a
suponer que las oportunidades fuera de ella también están igualmente
distribuidas, aun cuando las desigualdades de origen siguen operando más
allá de la escuela.

\section{Posición escolar y horizonte
educativo}\label{posiciuxf3n-escolar-y-horizonte-educativo}

Para los estudiantes escolares, la desigualdad de oportunidades se
materializa sobre todo en el acceso a la educación superior. Chile es un
caso paradigmático de masificación terciaria en América Latina
(\citeproc{ref-brunner_dinamicas_2016}{Brunner \& Ganga, 2016}): la
cobertura de pregrado pasó de cerca de 8\% en 1980 a 53\% en 2015
(\citeproc{ref-villalobos_percepciones_2022}{Villalobos et~al., 2022}).
La masificación, sin embargo, no eliminó la estratificación del acceso.
Según la encuesta CASEN 2009, la tasa de participación en educación
terciaria variaba entre 22,7\% en el quintil de menores ingresos y
67,8\% en el de mayores ingresos
(\citeproc{ref-brunner_lucha_2012}{Brunner, 2012}). Además, el sistema
se organiza como una pirámide de selectividad y reputación: Brunner
(\citeproc{ref-brunner_lucha_2012}{2012}) describe el ingreso a las
universidades chilenas como un proceso de doble selección, académica y
socioeconómica, en cuya cima predominan los egresados de colegios
particulares pagados, y el origen escolar del alumnado es uno de los
ejes que ordenan el sistema universitario
(\citeproc{ref-munoz_taxonomia_2013}{Muñoz \& Blanco, 2013}). Brunner
resume esta dinámica señalando que las universidades donde se forman las
élites reproducen ``bajo un velo meritocrático el privilegio de los
herederos'', al tiempo que ofrecen un canal de ascenso a quienes superan
las barreras selectivas.

El uso del mérito como cobertura del privilegio se ha observado también
desde dentro de las instituciones de élite. En un internado de élite
estadounidense, los estudiantes atribuyen su éxito a su trabajo duro,
aunque la observación muestra que rara vez trabajan intensamente y que
marginan a quienes sí lo hacen; la retórica meritocrática opera allí
como una cobertura de ventajas heredadas, en un mundo más abierto pero
más desigual (\citeproc{ref-khan_saying_2013}{Khan \& Jerolmack, 2013}).
Los estudiantes de Sciences Po y Oxford reconocen que viven en un mundo
desigual y que sus trayectorias están condicionadas por esa desigualdad,
pero al mismo tiempo creen firmemente que merecen el lugar que ocupan, y
recurren a la idea de responsabilidad para justificar su privilegio
(\citeproc{ref-allouch_ethique_2016}{Allouch et~al., 2016}). Los
estudiantes de Harvard y Stanford, por su parte, construyen fronteras
simbólicas entre universidades de élite que justifican su propia
posición (\citeproc{ref-binder_symbolically_2019}{Binder \& Abel,
2019}). Quienes acceden a las posiciones aventajadas tienden, así, a
interpretar su trayectoria como merecida.

En Chile, la evidencia sobre cómo los estudiantes perciben la
desigualdad en el acceso a la educación superior proviene de quienes ya
están dentro. Villalobos et al.
(\citeproc{ref-villalobos_percepciones_2022}{2022}) encuestaron a 1.773
estudiantes de ocho carreras de universidades de élite. El principal eje
de diferenciación de sus percepciones no fue el nivel socioeconómico,
sino el tipo de universidad: los estudiantes de universidades
tradicionales eran más escépticos de la meritocracia y de las
explicaciones individualistas de la desigualdad que sus pares de
universidades privadas elitistas. Los autores observan además
percepciones ``ortogonales'': los estudiantes explican la desigualdad
general por factores estructurales, pero atribuyen su propio ingreso a
sus cualidades personales.

En la escuela, la posición de cada estudiante se expresa ante todo en
sus calificaciones, que operan como marcadores institucionales de valor
y hacen visible el mérito de manera inmediata
(\citeproc{ref-laffert_changes}{Laffert et~al., 2026}). Al presentar los
resultados como ganados, los ideales meritocráticos alientan a los
``ganadores'' a ver su posición como merecida
(\citeproc{ref-laffert_changes}{Laffert et~al., 2026}). Desde una
perspectiva de \emph{interés propio}, quienes ocupan posiciones
ventajosas o son reconocidos por el sistema escolar tendrían razones
para percibir que las oportunidades están abiertas a todos, porque esa
percepción legitima su lugar (\citeproc{ref-baer_education_1982}{Baer \&
Lambert, 1982}; \citeproc{ref-dickinson_adolescent_1990}{Dickinson,
1990}; \citeproc{ref-duru-bellat_whos_2012}{Duru-Bellat \& Tenret,
2012}). Del mismo modo, en un sistema en que el acceso a la universidad
es selectivo y estratificado, quienes se proyectan en ella ocupan una
posición anticipada de ganadores del sistema.

\section{Trayectoria escolar: internalización e
ilustración}\label{trayectoria-escolar-internalizaciuxf3n-e-ilustraciuxf3n}

Una pregunta distinta es cómo cambia la percepción de igualdad de
oportunidades a medida que los estudiantes avanzan en la escuela y se
acercan a la transición a la educación superior. La hipótesis de
\emph{socialización} plantea que la educación transmite la ideología
dominante según el tiempo que se pasa en ella; la de \emph{ilustración},
que fomenta una mirada crítica sobre la desigualdad
(\citeproc{ref-baer_education_1982}{Baer \& Lambert, 1982}). Como ambos
procesos son, en sentido amplio, formas de socialización escolar, en
este estudio se denomina \emph{internalización} al primero, para
destacar que lo que se espera es la incorporación de la ideología
dominante. Los estudios sobre socialización universitaria muestran que
los contextos educativos pueden moldear actitudes sociopolíticas de
forma duradera (\citeproc{ref-alwin_political_1991}{Alwin et~al.,
1991}), y que la exposición a la institución modifica las percepciones:
entre estudiantes de élite chilenos, quienes llevan más tiempo en su
carrera atribuyen menos su ingreso a sus atributos personales
(\citeproc{ref-villalobos_percepciones_2022}{Villalobos et~al., 2022}).

En la etapa escolar, la evidencia es escasa y ambigua. Entre escolares
chilenos, los estudiantes de media usan la escala de percepción de
meritocracia de manera distinta a los de básica al evaluar si el
esfuerzo es recompensado, lo que se interpreta como un marco más crítico
entre los mayores (\citeproc{ref-laffert_changes}{Laffert et~al.,
2026}). Es plausible que, a medida que la selección universitaria se
vuelve más próxima y visible, los estudiantes perciban con mayor
claridad que no todos parten desde el mismo lugar.

\section{Hipótesis}\label{hipuxf3tesis}

A partir de los antecedentes se plantean cinco hipótesis, resumidas en
la Tabla~\ref{tbl-hipotesis}. Dos de ellas, referidas a las dimensiones
normativas de la creencia meritocrática, tienen un carácter exploratorio
y no anticipan una dirección.

\textbf{Transferencia.} Si la escuela es la primera institución en que
los estudiantes experimentan una distribución de recompensas que se
presenta como igualitaria y basada en el mérito, la evaluación de que
esa distribución es justa debería extenderse a la percepción de que la
sociedad ofrece las mismas oportunidades a todos
(\citeproc{ref-batruch_belief_2023}{Batruch et~al., 2023};
\citeproc{ref-mijs_stratified_2016}{Mijs, 2016a}). La hipótesis se
refiere a las percepciones de los estudiantes y no supone que la escuela
sea efectivamente justa o injusta.

\begin{quote}
\textbf{H1:} A mayor percepción de meritocracia en la escuela, mayor
percepción de igualdad de oportunidades en la sociedad.
\end{quote}

\textbf{Posición en la jerarquía escolar.} Quienes son reconocidos por
el sistema escolar tendrían razones para percibir que las oportunidades
están abiertas a todos, porque esa percepción legitima su lugar
(\citeproc{ref-baer_education_1982}{Baer \& Lambert, 1982};
\citeproc{ref-duru-bellat_whos_2012}{Duru-Bellat \& Tenret, 2012};
\citeproc{ref-laffert_changes}{Laffert et~al., 2026}).

\begin{quote}
\textbf{H2:} Los estudiantes con mejor rendimiento académico perciben
más igualdad de oportunidades.
\end{quote}

\textbf{Horizonte educativo.} Quienes se proyectan en la educación
superior, y en particular en la universidad, ocupan una posición
anticipada de ``ganadores'' del sistema. Como los estudiantes que ya
accedieron a instituciones selectivas, tenderían a interpretar su
trayectoria esperada como resultado de su esfuerzo y a percibir que las
oportunidades están abiertas
(\citeproc{ref-allouch_ethique_2016}{Allouch et~al., 2016};
\citeproc{ref-villalobos_percepciones_2022}{Villalobos et~al., 2022}).

\begin{quote}
\textbf{H3:} Los estudiantes que esperan completar una carrera
universitaria perciben más igualdad de oportunidades que quienes esperan
completar solo la educación media o una carrera técnica.
\end{quote}

\textbf{Internalización e ilustración.} La evidencia no permite
anticipar si el avance en la escolaridad refuerza o cuestiona la
percepción de igualdad de oportunidades
(\citeproc{ref-baer_education_1982}{Baer \& Lambert, 1982};
\citeproc{ref-duru-bellat_whos_2012}{Duru-Bellat \& Tenret, 2012};
\citeproc{ref-laffert_changes}{Laffert et~al., 2026}), por lo que se
plantea una hipótesis bilateral.

\begin{quote}
\textbf{H4:} La percepción de igualdad de oportunidades cambia con el
avance en la trayectoria escolar. Si predomina la internalización,
aumenta; si predomina la ilustración, disminuye.
\end{quote}

Para distinguir entre ambas explicaciones se examinará
complementariamente el cambio en la percepción de privilegio. La
ilustración implica que, junto con la menor percepción de igualdad de
oportunidades, aumenta la percepción de que el éxito depende de tener
padres ricos o buenos contactos.

\textbf{Creencias meritocráticas.} Se espera que las dimensiones
perceptivas de la creencia meritocrática se relacionen en direcciones
opuestas con la percepción de igualdad de oportunidades: la percepción
de meritocracia legitimaría las diferencias existentes y la percepción
de privilegio reconoce el peso del origen
(\citeproc{ref-castillo_meritocracia_2019}{Castillo et~al., 2019};
\citeproc{ref-laffert_changes}{Laffert et~al., 2026};
\citeproc{ref-mijs2016}{Mijs, 2016b}).

\begin{quote}
\textbf{H5a:} A mayor percepción de meritocracia, mayor percepción de
igualdad de oportunidades.
\end{quote}

\begin{quote}
\textbf{H5b:} A mayor percepción de privilegio, menor percepción de
igualdad de oportunidades.
\end{quote}

\textbf{Dimensiones normativas (exploratorias).} Para la preferencia
meritocrática y la aceptación del privilegio, los antecedentes apuntan
en direcciones contrapuestas: pueden operar como un estándar crítico
frente a la realidad o como justificaciones de la desigualdad
(\citeproc{ref-castillo_meritocracia_2019}{Castillo et~al., 2019};
\citeproc{ref-laffert_changes}{Laffert et~al., 2026}). Por ello se
examina su relación con la percepción de igualdad de oportunidades sin
anticipar su sentido.

\begin{quote}
\textbf{H5c (exploratoria):} La preferencia meritocrática se asocia con
la percepción de igualdad de oportunidades.
\end{quote}

\begin{quote}
\textbf{H5d (exploratoria):} La aceptación del privilegio se asocia con
la percepción de igualdad de oportunidades.
\end{quote}

\begin{longtable}[]{@{}
  >{\raggedright\arraybackslash}p{(\linewidth - 6\tabcolsep) * \real{0.2500}}
  >{\raggedright\arraybackslash}p{(\linewidth - 6\tabcolsep) * \real{0.2500}}
  >{\raggedright\arraybackslash}p{(\linewidth - 6\tabcolsep) * \real{0.2500}}
  >{\raggedright\arraybackslash}p{(\linewidth - 6\tabcolsep) * \real{0.2500}}@{}}
\caption{Resumen de hipótesis}\label{tbl-hipotesis}\tabularnewline
\toprule\noalign{}
\begin{minipage}[b]{\linewidth}\raggedright
Hipótesis
\end{minipage} & \begin{minipage}[b]{\linewidth}\raggedright
Mecanismo
\end{minipage} & \begin{minipage}[b]{\linewidth}\raggedright
Predictor
\end{minipage} & \begin{minipage}[b]{\linewidth}\raggedright
Efecto esperado
\end{minipage} \\
\midrule\noalign{}
\endfirsthead
\toprule\noalign{}
\begin{minipage}[b]{\linewidth}\raggedright
Hipótesis
\end{minipage} & \begin{minipage}[b]{\linewidth}\raggedright
Mecanismo
\end{minipage} & \begin{minipage}[b]{\linewidth}\raggedright
Predictor
\end{minipage} & \begin{minipage}[b]{\linewidth}\raggedright
Efecto esperado
\end{minipage} \\
\midrule\noalign{}
\endhead
\bottomrule\noalign{}
\endlastfoot
H1 & Transferencia & Meritocracia escolar percibida & Positivo \\
H2 & Interés propio & Rendimiento académico & Positivo \\
H3 & Horizonte educativo & Expectativa de completar la universidad &
Positivo \\
H4 & Internalización / ilustración & Ola (panel) & Positivo
(internalización) / negativo (ilustración) \\
H5a & Legitimación & Percepción de meritocracia & Positivo \\
H5b & Reconocimiento del origen & Percepción de privilegio & Negativo \\
H5c & Exploratoria & Preferencia meritocrática & Sin dirección
anticipada \\
H5d & Exploratoria & Aceptación del privilegio & Sin dirección
anticipada \\
\end{longtable}

\bookmarksetup{startatroot}

\chapter{Metodología}\label{metodologuxeda}

\section{Datos}\label{datos}

Se utilizan datos de la Encuesta de Percepciones de Meritocracia y
Formación Ciudadana, levantada en el marco del panel Educación y
Meritocracia (EDUMER), proyecto FONDECYT N.º 1210847
(\citeproc{ref-edumer_manual}{Proyecto Educación y Meritocracia, s.f.}).
Su objetivo es examinar las creencias, actitudes y comportamientos de
estudiantes en edad escolar respecto de la meritocracia, la desigualdad
y la ciudadanía (\citeproc{ref-laffert_changes}{Laffert et~al., 2026}).
El protocolo del estudio fue aprobado por un comité de ética, y se
registró la autorización de un adulto para la participación de cada
estudiante (\citeproc{ref-edumer_manual}{Proyecto Educación y
Meritocracia, s.f.}).

La encuesta es un panel de dos olas. La población objetivo son
estudiantes de escuelas municipales y particulares subvencionadas de las
regiones Metropolitana y de Valparaíso. Se encuestó a dos cohortes:
estudiantes de 6.º básico y 1.º medio en la primera ola, que fueron
seguidos en 7.º básico y 2.º medio en la segunda. Participaron nueve
escuelas (\citeproc{ref-laffert_changes}{Laffert et~al., 2026}). La
muestra no es probabilística ni representativa por zona geográfica o
dependencia administrativa (\citeproc{ref-edumer_manual}{Proyecto
Educación y Meritocracia, s.f.}).

El cuestionario fue autoadministrado en línea (CAWI) y respondido por
los estudiantes en las salas de computación de sus escuelas, con una
duración aproximada de 20 minutos. La primera ola se aplicó entre
octubre de 2023 y mayo de 2024, y la segunda entre agosto y noviembre de
2024 (\citeproc{ref-edumer_manual}{Proyecto Educación y Meritocracia,
s.f.}). La primera ola incluye 902 estudiantes y la segunda 674. De
ellos, 639 participan en ambas olas, 263 solo en la primera y 35 solo en
la segunda, lo que implica una atrición de 29,2\%
(\citeproc{ref-edumer_manual}{Proyecto Educación y Meritocracia, s.f.}).

Los análisis utilizan las observaciones de ambas olas, que suman 1.576
observaciones de 937 estudiantes. Los modelos se estiman con los casos
completos en todas sus variables, lo que deja 1.461 observaciones de 878
estudiantes, de los cuales 583 aportan observaciones en ambas olas.

La primera ola incluyó, al inicio del cuestionario, un módulo
experimental en el que los estudiantes fueron asignados aleatoriamente a
leer información sobre las diferencias de acceso a la universidad según
el tipo de colegio o un texto equivalente sobre la práctica deportiva
(\citeproc{ref-castillo_experimento_2024}{Castillo, Carrasco, et~al.,
2024}). Este módulo no se analiza en el presente estudio, y los modelos
no controlan por la asignación: una estimación preliminar con la primera
ola no encontró diferencias significativas en la percepción de igualdad
de oportunidades entre los estudiantes que recibieron la información y
el grupo control (razón de odds = 1,15; p = 0,28).

\section{Operacionalización}\label{operacionalizaciuxf3n}

\subsection{Variables}\label{variables}

La variable dependiente es la percepción de igualdad de oportunidades,
medida con el ítem ``En Chile, todas las personas tienen las mismas
oportunidades para salir adelante'' (\texttt{p1\_9}), tomado de los
cuestionarios de la Agencia de Calidad de la Educación. Dado que se
trata de un ítem de cuatro categorías ordenadas, se analiza como
variable ordinal; como robustez, se trata también como escala numérica
de 1 a 4 y como variable dicotómica que distingue a quienes están de
acuerdo o muy de acuerdo de quienes están en desacuerdo o muy en
desacuerdo. El ítem contiguo del cuestionario, ``En Chile, todas las
personas obtienen lo que merecen'' (\texttt{p1\_10}), no se combina con
la variable dependiente, porque se refiere a la regla que vincula el
mérito con la recompensa y ha sido utilizado como indicador de
percepción de meritocracia
(\citeproc{ref-castillo_meritocracia_2019}{Castillo et~al., 2019}).

La percepción de meritocracia escolar (H1) se mide con el ítem ``En esta
escuela, los/as estudiantes obtienen las notas que merecen''
(\texttt{p2\_3}), de modo que valores más altos indican una mayor
percepción de que la escuela distribuye las notas según el mérito. El
cuestionario incluye además dos ítems sobre si en la escuela se obtienen
buenas notas por el esfuerzo (\texttt{p2\_1}) o por la inteligencia
(\texttt{p2\_2}). Se optó por el ítem de merecimiento porque refiere
directamente a la evaluación de si la distribución de las notas es
merecida, que es lo que H1 supone que se extiende a la sociedad. El
rendimiento académico (H2) corresponde al promedio de notas del año
anterior (\texttt{p4}), en la escala chilena de 1,0 a 7,0. El horizonte
educativo (H3) se mide con la pregunta ``¿Cuál de los siguientes niveles
educacionales crees que vas a completar?'' (\texttt{p25}), dicotomizada
entre quienes esperan completar una carrera universitaria y quienes
esperan completar la educación media o una carrera en un instituto
profesional o centro de formación técnica. La trayectoria escolar (H4)
se mide con la ola de la encuesta, que distingue la primera de la
segunda medición de cada estudiante.

Las creencias meritocráticas (H5) se miden con la escala de percepciones
y preferencias sobre la meritocracia de Castillo et al.
(\citeproc{ref-castillo_multidimensional_2023}{2023}), que mide cada
dimensión con dos ítems (\citeproc{ref-laffert_changes}{Laffert et~al.,
2026}). La percepción de meritocracia (H5a) se mide con ``En Chile, las
personas son recompensadas por sus esfuerzos'' (\texttt{p1\_1}) y ``En
Chile, las personas son recompensadas por su inteligencia y habilidad''
(\texttt{p1\_2}) (\emph{r} = 0,50); la percepción de privilegio (H5b),
con ``En Chile, a quienes tienen padres ricos les va mucho mejor en la
vida'' (\texttt{p1\_3}) y ``En Chile, quienes tienen buenos contactos
les va mejor en la vida'' (\texttt{p1\_4}) (\emph{r} = 0,60); la
preferencia meritocrática (H5c), con ``Quienes más se esfuerzan deberían
obtener mayores recompensas que quienes se esfuerzan menos''
(\texttt{p1\_5}) y ``Quienes poseen más talento deberían obtener mayores
recompensas que quienes poseen menos talento'' (\texttt{p1\_6})
(\emph{r} = 0,37); y la aceptación del privilegio (H5d), con ``Está bien
que quienes tienen padres ricos les vaya bien en la vida''
(\texttt{p1\_7}) y ``Está bien que quienes tienen buenos contactos les
vaya bien en la vida'' (\texttt{p1\_8}) (\emph{r} = 0,61). Cada
dimensión se construye como el promedio de sus dos ítems; entre
paréntesis se indica la correlación policórica entre ellos en la primera
ola. El ítem de meritocracia escolar y los de la escala de creencias
meritocráticas se responden en cuatro categorías, desde muy en
desacuerdo (1) hasta muy de acuerdo (4), de modo que todas estas medidas
varían entre 1 y 4. Promediar los dos ítems de cada dimensión se
justifica en la evidencia psicométrica disponible para esta escala en
Laffert et~al. (\citeproc{ref-laffert_changes}{2026}). El artículo
reporta un análisis factorial confirmatorio y encuentra que las
creencias meritocráticas de los estudiantes se representan mejor con
cuatro factores distintos, con un ajuste aceptable en la mayoría de los
índices. Las cargas factoriales estandarizadas de los ocho ítems son
significativas y de moderadas a fuertes (entre 0,55 y 0,88), lo que
indica que los dos ítems de cada dimensión reflejan un mismo factor
latente y que su promedio puede interpretarse como una medida de esa
dimensión. Además, esta estructura alcanza invarianza estricta entre
olas dentro de los estudiantes.

\subsection{Controles}\label{controles}

Los modelos incluyen tres variables de control: género, país de
nacimiento y cohorte. Su objetivo es evitar que las relaciones entre los
predictores y la percepción de igualdad de oportunidades reflejen
diferencias en la composición de la muestra.

El género (\texttt{p20}) se incluye porque la percepción de meritocracia
varía según género en la población adulta
(\citeproc{ref-duru-bellat_whos_2012}{Duru-Bellat \& Tenret, 2012});
dado el reducido número de estudiantes que se identifica con otro
género, se distingue a los hombres de las mujeres y quienes se
identifican con otro género. El país de nacimiento del estudiante
(\texttt{p22\_1}), que distingue a quienes nacieron en Chile de quienes
nacieron en el extranjero, controla por una dimensión de desventaja
relevante en las escuelas estudiadas, en las que cerca de uno de cada
diez estudiantes nació fuera del país; como solo se preguntó en la
primera ola, se asigna a ambas observaciones de cada estudiante.

La cohorte distingue a los estudiantes que en la primera ola cursaban
6.º básico de quienes cursaban 1.º medio. Su inclusión responde a la
discusión sobre la trayectoria escolar. Si el avance en la escolaridad
fomenta una mirada más crítica sobre la desigualdad, como plantea la
hipótesis de ilustración (\citeproc{ref-baer_education_1982}{Baer \&
Lambert, 1982}), los estudiantes de media, con varios años más de
escolaridad, deberían percibir menos igualdad de oportunidades que los
de básica. Laffert et al. (\citeproc{ref-laffert_changes}{2026})
encuentran en este mismo panel que los estudiantes de media aplican un
marco más crítico al evaluar si el esfuerzo es recompensado. La ola y la
cohorte capturan así dos escalas de la trayectoria escolar: la ola, el
cambio de un año dentro de cada estudiante (H4); la cohorte, la
diferencia entre estudiantes separados por varios años de escolaridad.
Incluir la cohorte evita que esa diferencia se confunda con el cambio
entre olas, y su coeficiente se interpreta como evidencia complementaria
de H4. Como compara a dos grupos distintos de estudiantes, este
coeficiente puede reflejar también otras diferencias entre ellos, por lo
que no reemplaza al contraste dentro de los estudiantes.

El estudio no incluye una medida del nivel socioeconómico. El
cuestionario de estudiantes indaga la educación de la madre y del padre,
pero cerca de un cuarto de las observaciones no reporta la de la madre,
y la información de los apoderados solo puede vincularse a una fracción
minoritaria de los estudiantes. Se optó por no utilizar indicadores con
esa cobertura en lugar de introducir una medición poco confiable.

\subsection{Modelos (notaciones genneradas cn IA; pediennte de
confirmacion)}\label{modelos-notaciones-genneradas-cn-ia-pediennte-de-confirmacion}

Las hipótesis se contrastan con un modelo de regresión logística ordinal
que utiliza las observaciones de ambas olas. Como cada estudiante es
observado hasta dos veces, se estima un modelo multinivel de dos
niveles, en que las observaciones de cada ola (nivel 1) están anidadas
en estudiantes (nivel 2), con un intercepto aleatorio por estudiante.
Para el estudiante \(i\) en la ola \(t\):

\begin{equation}\protect\phantomsection\label{eq-modelo}{
\begin{aligned}
\text{logit}\,P(Y_{it} \le k) = \tau_k - \big(&\beta_1 \text{MeritEsc}_{it} + \beta_2 \text{Notas}_{it} + \beta_3 \text{Univ}_{it} + \beta_4 \text{Ola}_t \\
&+ \beta_5 \text{PercMerit}_{it} + \beta_6 \text{PercPriv}_{it} + \beta_7 \text{PrefMerit}_{it} + \beta_8 \text{AcepPriv}_{it} \\
&+ \boldsymbol{\delta}'\mathbf{X}_{it} + u_i\big)
\end{aligned}
}\end{equation}

donde \(\text{MeritEsc}_{it}\) es la percepción de meritocracia escolar,
\(\text{Notas}_{it}\) el promedio de notas, \(\text{Univ}_{it}\) la
expectativa de completar la universidad, \(\text{Ola}_t\) un indicador
de la segunda ola, \(\text{PercMerit}_{it}\), \(\text{PercPriv}_{it}\),
\(\text{PrefMerit}_{it}\) y \(\text{AcepPriv}_{it}\) las cuatro
dimensiones de la creencia meritocrática, \(\mathbf{X}_{it}\) el vector
de controles (género, país de nacimiento y cohorte) y
\(u_i \sim N(0, \sigma_u^2)\) el intercepto aleatorio del estudiante,
que recoge las diferencias estables entre estudiantes. Las variables
continuas se centran en su media general. Cada hipótesis corresponde a
un coeficiente: \(\beta_1 > 0\) (H1), \(\beta_2 > 0\) (H2),
\(\beta_3 > 0\) (H3), \(\beta_4 \neq 0\) (H4), cuyo signo permite
distinguir entre internalización (\(\beta_4 > 0\)) e ilustración
(\(\beta_4 < 0\)), \(\beta_5 > 0\) (H5a), \(\beta_6 < 0\) (H5b) y, de
manera exploratoria, \(\beta_7 \neq 0\) (H5c) y \(\beta_8 \neq 0\)
(H5d).

\section{Estrategia de análisis}\label{estrategia-de-anuxe1lisis}

El análisis se organiza en tres etapas.

\textbf{Descripción.} Se describirá la distribución de la percepción de
igualdad de oportunidades en cada ola y cohorte, y su correlación con la
percepción de meritocracia escolar y con las cuatro dimensiones de la
creencia meritocrática. Este paso permite evaluar empíricamente si los
estudiantes distinguen la igualdad de oportunidades de la meritocracia,
como sostiene el marco conceptual.

\textbf{Modelo panel (H1 a H5).} Las hipótesis se contrastarán con el
modelo descrito (Ecuación~\ref{eq-modelo}). Se reportarán los
coeficientes como razones de odds y la proporción de varianza explicada
por los efectos fijos y por el intercepto del estudiante (R² marginal y
condicional). Para distinguir la variación dentro de cada estudiante de
las diferencias estables entre estudiantes, los predictores que cambian
entre olas se descompondrán en su desviación respecto de la media del
estudiante y en dicha media. El componente dentro del estudiante indica
si los cambios de un mismo estudiante entre olas se asocian con cambios
en su percepción de igualdad de oportunidades, y no está afectado por
características estables no observadas.

\textbf{Robustez.} Se probarán modelos lineales sobre la escala de 1 a 4
y modelos logísticos sobre la versión dicotómica de la variable
dependiente; tambien se estimaran con difrerentes especificaciones de
los controles y covariables.

\bookmarksetup{startatroot}

\chapter{Resultados}\label{resultados}

\section{Análisis descriptivo}\label{anuxe1lisis-descriptivo}

\subsection{Percepción de igualdad de
oportunidades}\label{percepciuxf3n-de-igualdad-de-oportunidades}

Como la variable dependiente es un ítem ordinal de cuatro categorías, se
describe mediante su distribución completa y no mediante un promedio. La
Tabla~\ref{tbl-p19-ola} presenta la distribución por ola y la
Figura~\ref{fig-p19-distribucion} la desagrega por cohorte.

En la primera ola, el 39,3\% de los estudiantes está de acuerdo o muy de
acuerdo con que en Chile todas las personas tienen las mismas
oportunidades para salir adelante, mientras que el 60,7\% está en
desacuerdo o muy en desacuerdo. La categoría más frecuente es ``en
desacuerdo'' (36,2\%), que corresponde también a la mediana. La
percepción de igualdad de oportunidades es, por tanto, minoritaria entre
los estudiantes, aunque no marginal: cerca de cuatro de cada diez la
sostienen.

{

\begin{longtable}[]{@{}
  >{\raggedright\arraybackslash}p{(\linewidth - 4\tabcolsep) * \real{0.7714}}
  >{\raggedleft\arraybackslash}p{(\linewidth - 4\tabcolsep) * \real{0.1143}}
  >{\raggedleft\arraybackslash}p{(\linewidth - 4\tabcolsep) * \real{0.1143}}@{}}

\caption{\label{tbl-p19-ola}Percepción de igualdad de oportunidades
según ola. Frecuencias y porcentajes.}

\tabularnewline

\toprule\noalign{}
\begin{minipage}[b]{\linewidth}\raggedright
En Chile, todas las personas tienen las mismas oportunidades para salir
adelante
\end{minipage} & \begin{minipage}[b]{\linewidth}\raggedleft
Ola 1
\end{minipage} & \begin{minipage}[b]{\linewidth}\raggedleft
Ola 2
\end{minipage} \\
\midrule\noalign{}
\endhead
\bottomrule\noalign{}
\endlastfoot
Muy en desacuerdo & 220 (24,5\%) & 197 (29,2\%) \\
En desacuerdo & 325 (36,2\%) & 228 (33,8\%) \\
De acuerdo & 160 (17,8\%) & 141 (20,9\%) \\
Muy de acuerdo & 193 (21,5\%) & 108 (16,0\%) \\
Total & 898 & 674 \\

\end{longtable}

}

Las diferencias entre cohortes son marcadas. En la primera ola, el
52,9\% de los estudiantes de básica está de acuerdo o muy de acuerdo,
frente al 25,8\% de los de media. En básica, la categoría más frecuente
es ``muy de acuerdo'' (34,4\%), mientras que en media solo el 8,7\%
elige esa opción y el 74,2\% se ubica en el desacuerdo. La brecha se
mantiene en la segunda ola (48,3\% y 26,2\% de acuerdo,
respectivamente). Este patrón es consistente con el marco más crítico
que se ha observado entre los estudiantes de media al evaluar si el
esfuerzo es recompensado (\citeproc{ref-laffert_changes}{Laffert et~al.,
2026}), aunque la comparación entre cohortes debe interpretarse con
cautela, porque contrasta a grupos distintos de estudiantes y no el
cambio de los mismos estudiantes a lo largo del tiempo.

\begin{figure}

\caption{\label{fig-p19-distribucion}Percepción de igualdad de
oportunidades según ola y cohorte. Porcentajes; a la izquierda,
desacuerdo; a la derecha, acuerdo.}

\centering{

\pandocbounded{\includegraphics[keepaspectratio]{04-resultados_files/figure-pdf/fig-p19-distribucion-1.pdf}}

}

\end{figure}%

Entre olas, la percepción de igualdad de oportunidades tiende a
disminuir. El porcentaje de acuerdo pasa de 39,3\% a 36,9\%, y la
proporción que está muy en desacuerdo aumenta de 24,5\% a 29,2\%. Como
la composición de ambas olas difiere por la atrición, el cambio se
evalúa con los 635 estudiantes que responden en las dos mediciones
(Figura~\ref{fig-p19-transiciones}). De ellos, el 44,9\% mantiene la
misma respuesta, el 33,1\% se desplaza hacia el desacuerdo y el 22,0\%
hacia el acuerdo. La percepción es además relativamente inestable a
nivel individual: el 29,3\% de los estudiantes cambia de lado entre el
acuerdo y el desacuerdo de una ola a otra. Este descenso es consistente
con la hipótesis de ilustración (H5), que se evalúa formalmente en los
modelos multivariados.

\begin{figure}

\caption{\label{fig-p19-transiciones}Transiciones en la percepción de
igualdad de oportunidades entre la ola 1 y la ola 2. Estudiantes que
responden ambas olas; porcentajes sobre el total de esos estudiantes.
Los flujos tienen el color de la respuesta en la ola 1.}

\centering{

\pandocbounded{\includegraphics[keepaspectratio]{04-resultados_files/figure-pdf/fig-p19-transiciones-1.pdf}}

}

\end{figure}%

\subsection{Relación con las creencias
meritocráticas}\label{relaciuxf3n-con-las-creencias-meritocruxe1ticas}

Para evaluar si los estudiantes distinguen la igualdad de oportunidades
de las creencias meritocráticas, la Figura~\ref{fig-correlaciones}
presenta las correlaciones policóricas, adecuadas para ítems ordinales,
entre la variable dependiente y los ocho ítems de la escala de
percepciones y preferencias sobre la meritocracia, en la primera ola (n
entre 877 y 901 según el par de ítems).

\begin{figure}

\caption{\label{fig-correlaciones}Correlaciones policóricas entre la
percepción de igualdad de oportunidades y los ítems de la escala de
percepciones y preferencias sobre la meritocracia. Ola 1. Datos perdidos
tratados por pares (pairwise).}

\centering{

\pandocbounded{\includegraphics[keepaspectratio]{04-resultados_files/figure-pdf/fig-correlaciones-1.pdf}}

}

\end{figure}%

Las correlaciones de la igualdad de oportunidades con los demás ítems
son, en general, bajas o moderadas, y ninguna se acerca a la unidad. Las
más altas corresponden a la percepción de meritocracia: la igualdad de
oportunidades se asocia positivamente con percibir que en Chile se
recompensa el esfuerzo (r = 0,33) y la inteligencia (r = 0,27). Estas
correlaciones son menores que la que existe entre esos dos ítems entre
sí (r = 0,50), lo que indica que la percepción de igualdad de
oportunidades no es una medida más de la percepción de meritocracia,
sino una creencia relacionada pero distinta, como sostiene el marco
conceptual.

Las asociaciones con las dimensiones de privilegio siguen la dirección
esperada en las hipótesis. La percepción de privilegio se relaciona
negativamente con la igualdad de oportunidades (r = -0,27 para los
padres ricos y -0,26 para los contactos), mientras que la aceptación del
privilegio se relaciona positivamente (r = 0,18 y 0,23). La preferencia
meritocrática, en cambio, prácticamente no se asocia con la igualdad de
oportunidades (r = -0,09 y 0,00).

La matriz reproduce, además, la estructura de la escala de meritocracia:
los dos ítems de cada dimensión se correlacionan más entre sí (entre
0,37 y 0,61) que con los de otras dimensiones, lo que es consistente con
la estructura de cuatro factores validada en estos datos
(\citeproc{ref-laffert_changes}{Laffert et~al., 2026}). La percepción de
meritocracia y la de privilegio se relacionan de forma débil y
mayormente negativa, lo que muestra que los estudiantes no las tratan
como polos opuestos de una misma creencia. Las correlaciones entre los
ítems de cada dimensión son, además, similares a las que se desprenden
del análisis factorial confirmatorio que Laffert et al.
(\citeproc{ref-laffert_changes}{2026}) estimaron con estos mismos datos.

\bookmarksetup{startatroot}

\chapter{Discusión}\label{discusiuxf3n}

Pendiente.

\section{Limitaciones}\label{limitaciones}

Por lo pronto no hay forma consistente de operacionalizar el NSE, poca
varianza en educación de los padres, discordancia entre los niveles
educativos reportados por los estudiantes y los apoderados, cantidad de
libros como proxy de capital cultural tampoco en este caso reporta
efectos significativos en ninguna de sus especificaciones, ni en versión
ordinal, ni en más de veinticinco versus el resto, los factores de
meritocracia estarían entrando como índices ponderados; se están sumando
los errores, habría que considerar algún tipo de análisis de robustez,
me da la impresión de que la inclusión del experimento está interesante
más no del todo justificada, de cualquier forma sirve para controlar al
menos en la ola uno si la info que se entrega a los estudiantes tiene
algún efecto sobre el resto de respuestas de la encuesta; ``¿Conocer
cifras de desigualdad en el acceso a la universidad afecta la percepcion
de igualdad de oportunidades?'' se me ocurre que tal vez podría
imputarse el control a la ola dos, si los estudiantes en algún momento
se les entregó la información sobre desigualdad, podría ser que para el
cuestionario dos se hayan acordado y también tenga efecto sobre sus
respuestas.

\bookmarksetup{startatroot}

\chapter{Conclusiones}\label{conclusiones}

Síntesis y proyecciones futuras.

\bookmarksetup{startatroot}

\chapter{Referencias}\label{referencias}

\protect\phantomsection\label{refs}
\begin{CSLReferences}{1}{0}
\bibitem[\citeproctext]{ref-allouch_ethique_2016}
Allouch, A., Brown, P., Power, S., \& Tholen, G. (2016). L'éthique des
élites scolaires. Du mérite à la responsabilité chez les étudiants de
Sciences Po Paris et de l'Université d'Oxford. \emph{L'Année
sociologique}, \emph{66}(1), 193-224.
\url{https://doi.org/10.3917/anso.161.0193}

\bibitem[\citeproctext]{ref-alwin_political_1991}
Alwin, D. F., Cohen, R. L., \& Newcomb, T. M. (1991). \emph{Political
Attitudes over the Life Span: The Bennington Women after Fifty Years}.
University of Wisconsin Press.

\bibitem[\citeproctext]{ref-baer_education_1982}
Baer, D. E., \& Lambert, R. D. (1982). Education and Support for
Dominant Ideology. \emph{Canadian Review of Sociology / Revue canadienne
de sociologie}, \emph{19}(2), 173-195.
\url{https://doi.org/10.1111/j.1755-618X.1982.tb00859.x}

\bibitem[\citeproctext]{ref-batruch_belief_2023}
Batruch, A., Jetten, J., Van De Werfhorst, H., Darnon, C., \& Butera, F.
(2023). Belief in {School Meritocracy} and the {Legitimization} of
{Social} and {Income Inequality}. \emph{Social Psychological and
Personality Science}, \emph{14}(5), 621-635.
\url{https://doi.org/10.1177/19485506221111017}

\bibitem[\citeproctext]{ref-binder_symbolically_2019}
Binder, A. J., \& Abel, A. R. (2019). Symbolically Maintained
Inequality: How Harvard and Stanford Students Construct Boundaries among
Elite Universities. \emph{Sociology of Education}, \emph{92}(1), 41-58.
\url{https://doi.org/10.1177/0038040718821073}

\bibitem[\citeproctext]{ref-bourdieu_reproduction_1990}
Bourdieu, P., \& Passeron, J.-C. (1990). \emph{Reproduction in
Education, Society and Culture} (2.ª ed.). Sage.

\bibitem[\citeproctext]{ref-brunner_lucha_2012}
Brunner, J. J. (2012). La lucha por la educación de las elites: campo y
canales formativos. \emph{Revista UDP}, \emph{9}, 119-143.

\bibitem[\citeproctext]{ref-brunner_dinamicas_2016}
Brunner, J. J., \& Ganga, F. A. (2016). Dinámicas de transformación en
la educación superior latinoamericana: desafíos para la gobernanza.
\emph{Opción}, \emph{32}(80), 12-35.

\bibitem[\citeproctext]{ref-castillo_experimento_2024}
Castillo, J. C., Carrasco, K., Nuñez, M. F., \& Laffert, A. (2024).
\emph{Meritocracia en la etapa escolar: un experimento de encuesta}.
Presentación en el Congreso Chileno de Sociología, Santiago, 2-4 de
mayo. \url{https://github.com/educacion-meritocracia/experimento}

\bibitem[\citeproctext]{ref-castillo_multidimensional_2023}
Castillo, J. C., Iturra, J., Maldonado, L., Atria, J., \& Meneses, F.
(2023). A {Multidimensional Approach} for {Measuring Meritocratic
Beliefs}: {Advantages}, {Limitations} and {Alternatives} to the {ISSP
Social Inequality Survey}. \emph{International Journal of Sociology},
\emph{53}(6), 448-472.
\url{https://doi.org/10.1080/00207659.2023.2274712}

\bibitem[\citeproctext]{ref-castillo_percepcion_2012}
Castillo, J. C., Miranda, D., \& Carrasco, D. (2012). Percepción de
desigualdad económica en Chile: medición, diferencias y determinantes.
\emph{Psykhe}, \emph{21}(1), 99-114.

\bibitem[\citeproctext]{ref-castillo_socialization_2024}
Castillo, J. C., Salgado, M., Carrasco, K., \& Laffert, A. (2024). The
{Socialization} of {Meritocracy} and {Market Justice Preferences} at
{School}. \emph{Societies}, \emph{14}(11), 214.
\url{https://doi.org/10.3390/soc14110214}

\bibitem[\citeproctext]{ref-castillo_meritocracia_2019}
Castillo, J. C., Torres, A., Atria, J., \& Maldonado, L. (2019).
Meritocracia y Desigualdad Económica: {Percepciones}, Preferencias e
Implicancias. \emph{Revista Internacional de Sociología}, \emph{77}(1),
e117. \url{https://doi.org/10.3989/ris.2019.77.1.17.114}

\bibitem[\citeproctext]{ref-darnon_where_2018}
Darnon, C., Wiederkehr, V., Dompnier, B., \& Martinot, D. (2018).
{«{Where} There Is a Will, There Is a Way»}: {Belief} in School
Meritocracy and the Social-class Achievement Gap. \emph{British Journal
of Social Psychology}, \emph{57}(1), 250-262.
\url{https://doi.org/10.1111/bjso.12214}

\bibitem[\citeproctext]{ref-dickinson_adolescent_1990}
Dickinson, J. (1990). Adolescent Representations of Socio-economic
Status. \emph{British Journal of Developmental Psychology}, \emph{8}(4),
351-371. \url{https://doi.org/10.1111/j.2044-835X.1990.tb00850.x}

\bibitem[\citeproctext]{ref-duru-bellat_whos_2012}
Duru-Bellat, M., \& Tenret, E. (2012). Who's for Meritocracy?
{Individual} and Contextual Variations in the Faith. \emph{Comparative
Education Review}, \emph{56}(2), 223-247.
\url{https://doi.org/10.1086/661290}

\bibitem[\citeproctext]{ref-goudeau_hidden_2017}
Goudeau, S., \& Croizet, J.-C. (2017). Hidden Advantages and
Disadvantages of Social Class: {How} Classroom Settings Reproduce Social
Inequality by Staging Unfair Comparison. \emph{Psychological Science},
\emph{28}(2), 162-170. \url{https://doi.org/10.1177/0956797616676600}

\bibitem[\citeproctext]{ref-khan_saying_2013}
Khan, S., \& Jerolmack, C. (2013). Saying Meritocracy and Doing
Privilege. \emph{The Sociological Quarterly}, \emph{54}(1), 9-19.
\url{https://doi.org/10.1111/tsq.12008}

\bibitem[\citeproctext]{ref-laffert_changes}
Laffert, A., Castillo, J. C., Canales, R., Urzúa, T., \& Carrasco, K.
(2026). \emph{Stability and comparability of meritocratic beliefs in
school-age students: A measurement invariance approach across time and
cohorts}. SocArXiv. \url{https://doi.org/10.31235/osf.io/gvbqw_v1}

\bibitem[\citeproctext]{ref-mijs_stratified_2016}
Mijs, J. J. B. (2016a). Stratified Failure: {Educational} Stratification
and Students' Attributions of Their Mathematics Performance in 24
Countries. \emph{Sociology of Education}, \emph{89}(2), 137-153.
\url{https://doi.org/10.1177/0038040716636434}

\bibitem[\citeproctext]{ref-mijs2016}
Mijs, J. J. B. (2016b). \emph{The Unfulfillable Promise of Meritocracy:
Three Lessons and Their Implications for Justice in Education}.

\bibitem[\citeproctext]{ref-munoz_taxonomia_2013}
Muñoz, M., \& Blanco, C. (2013). Una taxonomía de las universidades
chilenas. \emph{Calidad en la Educación}, \emph{38}.

\bibitem[\citeproctext]{ref-edumer_manual}
Proyecto Educación y Meritocracia. (s.f.). \emph{Encuesta Panel
Educación y Meritocracia (EDUMER): manual metodológico}.
\url{https://educacion-meritocracia.github.io/encuesta-edumer/}

\bibitem[\citeproctext]{ref-resh2010}
Resh, N. (2010). Sense of Justice about Grades in School: Is It
Stratified like Academic Achievement? \emph{Social Psychology of
Education}, \emph{13}, 313-329.
\url{https://doi.org/10.1007/s11218-010-9117-z}

\bibitem[\citeproctext]{ref-resh_sense_2017}
Resh, N., \& Sabbagh, C. (2017). Sense of Justice in School and Civic
Behavior. \emph{Social Psychology of Education}, \emph{20}(2), 387-409.
\url{https://doi.org/10.1007/s11218-017-9375-0}

\bibitem[\citeproctext]{ref-villalobos_percepciones_2022}
Villalobos, C., Kuzmanic, D., Valenzuela, J. P., \& Quaresma, M. L.
(2022). Percepciones sociales de los estudiantes de educación superior
de élite: un estudio del caso chileno. \emph{Psykhe}, \emph{31}(2),
1-18. \url{https://doi.org/10.7764/psykhe.2020.22549}

\bibitem[\citeproctext]{ref-wiederkehr_belief_2015}
Wiederkehr, V., Bonnot, V., Krauth-Gruber, S., \& Darnon, C. (2015).
Belief in School Meritocracy as a System-Justifying Tool for Low Status
Students. \emph{Frontiers in Psychology}, \emph{6}.
\url{https://doi.org/10.3389/fpsyg.2015.01053}

\bibitem[\citeproctext]{ref-young1958}
Young, M. (1958). \emph{The Rise of the Meritocracy, 1870--2033}.
Thames; Hudson.

\end{CSLReferences}

\cleardoublepage
\phantomsection
\addcontentsline{toc}{part}{Apéndices}
\appendix

\chapter{Instrumentos de
levantamiento}\label{instrumentos-de-levantamiento}

Incluye cuestionarios, guias de entrevista, etc.

\chapter{Tablas adicionales}\label{tablas-adicionales}

Tablas extendidas, tests de robustez, etc.


\backmatter

% --- Cuerpo del libro (capítulos) en arábigos ---
\mainmatter
\pagestyle{scrheadings}

% --- (Opcional) Estilo del índice general ---
% \addtocontents{toc}{\protect\thispagestyle{plain}}


\end{document}
